\section{Execution granularity and local programs}\label{sec:C18:granularity}
One model decision per atomic tool call exposes every intermediate value to the main reasoning trace. This can be useful when each observation changes the plan, but wasteful when several operations form a locally deterministic subroutine. Code-based composition can keep temporary values inside an execution environment and return only the result needed by the outer policy. It changes execution granularity, not merely trace compression.

HyperTool uses an executable block that can call existing tools through their schemas and manipulate intermediate values locally~\cite{du2026hypertool}. Its method couples that interface with synthesized and verified training trajectories. The reported gains therefore involve both a new action abstraction and training for its use. A comparison that gives only one condition the specialized training cannot attribute the full result to block execution alone. The public paper now permits a method-level assessment beyond an abstract-only lead, while independent reproduction remains unestablished here.

The mathematical distinction is between choosing a sequence of atomic actions with model feedback at every step and choosing a program $p$ whose internal transition sequence is executed under a runtime. A program can express loops, branches, and dataflow compactly. Its correctness requires both valid tool calls and correct control flow. It can reduce model-call latency while making intermediate effects less visible to the outer policy. The tradeoff depends on how often new observations should trigger reconsideration.

A constructed example retrieves two coordinates and computes a distance. If geocoding and arithmetic are deterministic under the request, one block can pass values locally without asking the model to copy them between calls. If the first geocoding result is ambiguous, the block should expose that ambiguity or apply a declared rule rather than silently choosing a location. Grouping does not eliminate the need for checkpoints at semantically meaningful uncertainty boundaries.

Security enforcement must cover calls inside the block. Wrapping the outer executable tool while allowing unrestricted internal tool access can bypass per-tool policy. Every consequential internal invocation should preserve the principal, authority, operation identity, and audit trail. A returned concise result should retain references to the internal effect receipts. Context savings should not be purchased by deleting the evidence needed for recovery.

The proposed ablation compares atomic calls, an untrained block interface, and a block-trained policy under matched models and budgets. Measure model tokens, internal tool calls, critical-path latency, error recovery, and hidden effects. A reduction in outer calls is informative but incomplete. The relevant question is whether the new granularity preserves justified task behavior while reducing complete execution cost.

